
%(BEGIN_QUESTION)
% Copyright 2006, Tony R. Kuphaldt, released under the Creative Commons Attribution License (v 1.0)
% This means you may do almost anything with this work of mine, so long as you give me proper credit

Hvis underdimensjonering av en reguleringsventil fører til at man ikke oppnår ønsket strømningsrate, hvorfor ikke bare overdimensjonere alle reguleringsventiler for å sikre at ønskede strømningsrater kan oppnås? Hva ville være galt med å bruke denne filosofien for ventildimensjonering?

\underbar{file i01422}
%(END_QUESTION)





%(BEGIN_ANSWER)

Overdimensjonerte reguleringsventiler vil riktignok kunne oppnå de ønskede strømningsratene, men de gjør det umulig å {\it regulere} strømningen presist, spesielt ved lave strømningsrater.

\vskip 10pt

For å forstå hvorfor, får du dette ``tankeeksperimentet'': Se for deg at den lille håndventilen på kjøkkenkranen byttes ut med en ekstremt stor ventil, for eksempel på størrelse med en brannhydrantventil. Med denne absurd overdimensjonerte ventilen koblet til det samme rørsystemet i huset ditt, hvordan ville det være å prøve å regulere vannstrømmen i kjøkkenvasken?
 
%(END_ANSWER)





%(BEGIN_NOTES)

%INDEX% Final Control Elements, valve: undersizing

%(END_NOTES)

